\documentclass[12pt]{article}

\usepackage[T1]{fontenc}
\usepackage[utf8]{inputenc}
\usepackage[romanian]{babel}

\usepackage{amsmath, amssymb, amsthm}
\usepackage{mathrsfs}
\usepackage{geometry}
\geometry{a4paper, margin=2.5cm}

\title{Histograme și Variabile Aleatoare Discrete}
\author{}
\date{}

\theoremstyle{plain}
\newtheorem{definition}{Definiție}
\newtheorem{proposition}{Propoziție}
\newtheorem{example}{Exemplu}

\begin{document}

\maketitle

\section{Histograme}

\begin{definition}
Fie $X$ o variabilă aleatoare reală și fie $x_1, \dots, x_n$ un eșantion
independent extras din distribuția lui $X$.
\textbf{Histograma} eșantionului este funcția care asociază fiecărui interval
$B$ numărul (sau proporția) de observații care cad în $B$:


\[
H_n(B) := \frac{1}{n} \sum_{i=1}^n \mathbf{1}_{\{x_i \in B\}}.
\]


\end{definition}

Histograma este o aproximare empirică a distribuției lui $X$:


\[
H_n(B) \approx \mathbb{P}(X \in B).
\]



\section{Variabile aleatoare discrete}

\begin{definition}
O variabilă aleatoare $X : \Omega \to S$ se numește
\textbf{discretă} dacă imaginea sa este cel mult numărabilă:


\[
X(\Omega) = \{x_1, x_2, \dots\}.
\]


Distribuția lui $X$ este determinată de probabilitățile:


\[
p_k := \mathbb{P}(X = x_k).
\]


Funcția $p : S \to [0,1]$ definită prin $p(x_k) = p_k$ se numește
\textbf{funcție de masă a probabilității}.
\end{definition}

\section{Exemple fundamentale}

\subsection{Distribuția Bernoulli}

\begin{definition}
O variabilă aleatoare $X$ se numește \textbf{Bernoulli} cu parametru $p \in [0,1]$ dacă:


\[
\mathbb{P}(X = 1) = p, \qquad \mathbb{P}(X = 0) = 1 - p.
\]


Notăm $X \sim \mathrm{Bernoulli}(p)$.
\end{definition}

\subsection{Distribuția polinomială pe o mulțime finită}

\begin{definition}
Fie $S = \{x_1, \dots, x_k\}$ o mulțime finită și $p_1, \dots, p_k$ probabilități cu
$\sum_{i=1}^k p_i = 1$.
O variabilă aleatoare $X$ cu:


\[
\mathbb{P}(X = x_i) = p_i
\]


se numește \textbf{distribuție polinomială} pe $S$.
\end{definition}

\subsection{Distribuția geometrică}

\begin{definition}
O variabilă aleatoare $X$ se numește \textbf{geometrică} cu parametru $p \in (0,1]$ dacă:


\[
\mathbb{P}(X = k) = (1-p)^{k-1} p, \qquad k = 1,2,\dots
\]


Notăm $X \sim \mathrm{Geom}(p)$.
\end{definition}

\subsection{Distribuția binomială negativă}

\begin{definition}
Fie $r \in \mathbb{N}$ și $p \in (0,1]$.
O variabilă aleatoare $X$ se numește \textbf{binomială negativă} dacă:


\[
\mathbb{P}(X = k) = \binom{k-1}{r-1} p^r (1-p)^{k-r},
\qquad k = r, r+1, \dots
\]


Notăm $X \sim \mathrm{NegBin}(r,p)$.
\end{definition}

\subsection{Distribuția Poisson}

\begin{definition}
O variabilă aleatoare $X$ se numește \textbf{Poisson} cu parametru $\lambda > 0$ dacă:


\[
\mathbb{P}(X = k) = e^{-\lambda} \frac{\lambda^k}{k!},
\qquad k = 0,1,2,\dots
\]


Notăm $X \sim \mathrm{Poisson}(\lambda)$.
\end{definition}

\subsection{Distribuția hipergeometrică}

\begin{definition}
Fie o populație cu $N$ elemente, dintre care $M$ sunt de tip ``succes''.
Extragem fără înlocuire $n$ elemente.
Variabila aleatoare $X$ = numărul de succese are distribuție \textbf{hipergeometrică}:


\[
\mathbb{P}(X = k)
= \frac{\binom{M}{k} \binom{N-M}{n-k}}{\binom{N}{n}},
\qquad k = 0,1,\dots,n.
\]


Notăm $X \sim \mathrm{Hypergeom}(N,M,n)$.
\end{definition}

\subsection{Distribuția multinomială}

\begin{definition}
Fie $n \in \mathbb{N}$ și probabilități $p_1,\dots,p_k$ cu $\sum p_i = 1$.
Un vector aleator


\[
(X_1,\dots,X_k)
\]


are distribuție \textbf{multinomială} dacă:


\[
\mathbb{P}(X_1 = x_1, \dots, X_k = x_k)
= \frac{n!}{x_1! \cdots x_k!} p_1^{x_1} \cdots p_k^{x_k},
\]


unde $x_1 + \cdots + x_k = n$.
Notăm $(X_1,\dots,X_k) \sim \mathrm{Multinomial}(n; p_1,\dots,p_k)$.
\end{definition}


\end{document}
